\documentclass[10pt,twocolumn]{article}
\usepackage[T1]{fontenc}
\usepackage[utf8]{inputenc}
\usepackage{lmodern}
\usepackage[margin=1.8cm,top=2.4cm,bottom=2cm,columnsep=0.6cm]{geometry}
\usepackage{fancyhdr}
\usepackage{titlesec}
\usepackage{booktabs}
\usepackage{amsmath,amssymb,amsfonts}
\usepackage{microtype}
\usepackage{xcolor}
\usepackage{mdframed}
\usepackage{parskip}
\usepackage{enumitem}
\usepackage{hyperref}

\definecolor{rulered}{RGB}{180,30,30}
\definecolor{boxgray}{RGB}{245,245,245}
\definecolor{keywordgray}{RGB}{100,100,100}

\pagestyle{fancy}
\fancyhf{}
\fancyhead[L]{\textsc{\small Claude Mag}}
\fancyhead[C]{\small\textit{Machine Learning Research Digest}}
\fancyhead[R]{\small 2026-08-10}
\fancyfoot[C]{\small\thepage}
\renewcommand{\headrulewidth}{0.8pt}

\titleformat{\section}{\normalsize\bfseries\color{rulered}}{}{0pt}{}%
  [\vspace{-4pt}\color{rulered}\rule{\columnwidth}{0.4pt}]
\titlespacing{\section}{0pt}{8pt}{4pt}

\hypersetup{colorlinks=true,urlcolor=rulered,linkcolor=black}

\begin{document}
\setlength{\parindent}{0pt}
\setlength{\parskip}{4pt}

{\small\color{keywordgray}\textbf{Keywords:}
  offline reinforcement learning \textbullet{}
  diffusion policy \textbullet{}
  Q-value correction \textbullet{}
  distribution shift}\par\vspace{4pt}

{\LARGE\bfseries\raggedright
  Diffusion Policy with Behavioral Advantage Correction
  for Offline Reinforcement Learning}\par\vspace{4pt}

{\large\itshape\raggedright
  BAC-PE Uses the Behavior Policy's Q-Function to Simultaneously
  Tame Overestimation and Conservatism in Offline RL}\par\vspace{6pt}

{\small\textbf{By} Botao Dong, Longyang Huang, Ning Pang,
  Hongtian Chen}\par\vspace{2pt}

{\small\color{keywordgray}arXiv:2608.02332 \textbullet{} August 2026}
\par\vspace{2pt}\noindent{\color{rulered}\rule{\columnwidth}{0.6pt}}\vspace{6pt}

Offline reinforcement learning is caught between two failure modes: over-conservative
penalization of in-distribution actions and overestimation of out-of-distribution
actions. Behavioral Advantage Corrected Policy Evaluation (BAC-PE) resolves this
dilemma by adding the behavior policy's advantage function as a correction term to
the learned Q-function, with convergence guarantees and strong empirical results
across D4RL, Antmaze, and robotic manipulation benchmarks.

\begin{mdframed}[backgroundcolor=boxgray,linecolor=rulered,linewidth=1pt,
    innerleftmargin=6pt,innerrightmargin=6pt,
    innertopmargin=6pt,innerbottommargin=6pt]
\textbf{\small AT A GLANCE}\par\vspace{4pt}
\begin{description}[leftmargin=0pt,labelsep=4pt,itemsep=2pt,parsep=0pt,
    font=\small\bfseries]
  \item[Problem:] Distribution shift in offline RL causes Q-functions to either
    overestimate out-of-distribution actions or over-penalize in-distribution ones.
  \item[Why it matters:] Safe offline RL from logged data is essential for robotics
    and healthcare, where online interaction is expensive or dangerous.
  \item[Method:] Correct the learned Q-function with the behavior policy's advantage;
    represent both policies as diffusion models for multimodal action distributions.
  \item[Result:] Outperforms CQL, TD3+BC, and Diffusion-QL on D4RL locomotion and
    Antmaze; strong results on robotic manipulation.
  \item[Evidence:] Theoretical convergence bound on Q-estimation error; evaluated on
    D4RL, Antmaze, and Robosuite benchmarks.
\end{description}
\end{mdframed}

\section{The Big Picture}

Offline reinforcement learning promises to extract useful policies from logs of past
behavior---a critical capability for domains such as healthcare and autonomous driving
where online interaction is too costly or dangerous. The fundamental challenge is that
the learned policy will inevitably visit state-action pairs not well represented in
the fixed dataset. Standard Q-learning then propagates erroneous value estimates for
these out-of-distribution (OOD) actions, causing the policy to exploit them
catastrophically.

The field has responded with two families of corrections. \emph{Conservative} methods
(CQL, IQL) penalize OOD actions with explicit regularization, but the penalty is
necessarily blunt: it can suppress well-covered in-distribution actions alongside
genuinely dangerous OOD ones. \emph{Implicit constraint} methods (TD3+BC, BEAR)
restrict the policy to stay close to the behavior distribution, but tight constraints
sacrifice the ability to improve over a mediocre behavior policy. BAC-PE takes a
fundamentally different approach: rather than penalizing OOD actions by their
distance from the behavior distribution, it \emph{corrects} Q-values using
information the behavior policy actually provides.

\section{Why Existing Approaches Fall Short}

The behavior policy $\mu$ generates all data in the offline dataset. Its Q-function
$Q^\mu$ is estimable from the data without distribution shift, because every action
in the dataset was taken by $\mu$ on the state distribution that $\mu$ itself induced.
This is the key observation that existing methods leave on the table.

\textbf{Conservative Q-Learning (CQL)} adds a penalty $-\mathbb{E}_{a \sim \pi}[Q]$
to push Q-values down for policy actions and up for data actions. But it does so
uniformly across all policy actions, regardless of whether they are genuinely OOD or
merely low-density actions that $\mu$ took rarely but which are still valid.

\textbf{TD3+BC} adds a behavioral cloning regularizer $\|a - a_\mu\|^2$ to the
policy gradient, preventing large deviations from the logged actions but also
preventing the policy from taking shortcuts the behavior policy missed.

Neither method uses $Q^\mu$ directly. BAC-PE does.

\section{Core Mechanism}

\textbf{Behavioral Advantage.}
The behavioral advantage of action $a$ in state $s$ is:
\[
A^\mu(s, a) = Q^\mu(s, a) - V^\mu(s),
\]
where $V^\mu(s) = \mathbb{E}_{a \sim \mu}[Q^\mu(s,a)]$.
A positive behavioral advantage means $a$ is \emph{better than average} under $\mu$;
a negative behavioral advantage means $a$ is OOD or poor.

\textbf{BAC-PE Correction.}
The corrected Q-function for the learned policy $\pi$ is:
\[
\hat{Q}^\pi(s, a) = Q^\pi(s, a) + \alpha\,A^\mu(s, a).
\]
The intuition:
\begin{itemize}[itemsep=2pt,parsep=0pt]
  \item If $a$ was frequently taken and good under $\mu$ (large $A^\mu > 0$),
    $\hat{Q}^\pi$ is boosted, offsetting excessive conservatism.
  \item If $a$ was rare or poor under $\mu$ (small or negative $A^\mu$),
    $\hat{Q}^\pi$ is suppressed, offsetting overestimation.
\end{itemize}
No uniform penalty is applied; the correction is adaptive, action-by-action.

\textbf{Diffusion Policy.}
Both $\mu$ and $\pi$ are represented as \emph{diffusion models} over the action space:
\[
\pi_\theta(a|s) \propto \int p_{\theta,T}(a_T)\prod_{t=1}^T
p_{\theta,t}(a_{t-1}|a_t,\,s)\,\mathrm{d}a_{1:T},
\]
where $a_T \sim \mathcal{N}(0, I)$ and $p_{\theta,t}$ is a learned denoising step.
Diffusion policies capture multimodal action distributions (e.g.\ multiple valid
grasps for a manipulation task) that Gaussian or deterministic policies cannot.

\section{Technical Formulation}

Let $\mathcal{D} = \{(s_i, a_i, r_i, s'_i)\}$ be the offline dataset collected under
$\mu$. BAC-PE estimates $Q^\mu$ and $V^\mu$ via standard TD with dataset transitions:
\[
Q^\mu_{\phi} \leftarrow \arg\min_\phi \mathbb{E}_{\mathcal{D}}\!\left[
  \bigl(Q^\mu_\phi(s,a) - r - \gamma\,V^\mu_\phi(s')\bigr)^2\right],
\]
then learns $Q^\pi$ with the corrected target:
\[
\hat{Q}^\pi_\psi \leftarrow \arg\min_\psi \mathbb{E}_{\mathcal{D}}\!\left[
  \Bigl(\hat{Q}^\pi_\psi(s,a) - r - \gamma\,\hat{Q}^\pi_\psi(s', a'_{s'})\Bigr)^2
  \right],
\]
where $a'_{s'} \sim \pi_\theta(\cdot|s')$.

The policy is optimized to maximize $\hat{Q}^\pi$ subject to a KL constraint on
deviation from $\mu$:
\[
\mathcal{L}(\theta) = -\mathbb{E}_{s, a \sim \pi_\theta}\!\left[\hat{Q}^\pi(s,a)\right]
+ \lambda\,D_{\mathrm{KL}}(\pi_\theta(\cdot|s)\;\|\;\mu(\cdot|s)).
\]

\section{Learning or Inference Procedure}

\textbf{Stage 1: Behavior policy fitting.} A diffusion model is fit to the offline
dataset actions conditioned on states: this is standard behavioral cloning with a
diffusion denoiser.

\textbf{Stage 2: Q-function estimation.} $Q^\mu_\phi$ and $V^\mu_\phi$ are trained
via TD using the offline dataset.

\textbf{Stage 3: Policy optimization.} The diffusion policy $\pi_\theta$ is optimized
alternately updating the denoising network (to maximize $\hat{Q}^\pi$) and the
Q-function $Q^\pi_\psi$ (to provide targets for the next policy update).

\textbf{Inference.} At test time, the diffusion policy generates actions by running $T$
denoising steps from Gaussian noise conditioned on the current state.

\section{What the Guarantee Says}

\textbf{Theorem (informal).} Under standard concentrability and function approximation
assumptions, the BAC-PE Q-function $Q^\pi_k$ at iteration $k$ satisfies:
\[
\|Q^\pi_k - Q^{\pi^*}\|_\infty \;\leq\; \varepsilon_k + C\,\delta_\mu,
\]
where $\varepsilon_k \to 0$ as $k \to \infty$, $C$ is a problem constant, and
$\delta_\mu = \sup_{s,a}(\pi(a|s)/\mu(a|s))$ is the concentrability coefficient.
Actions well-covered by $\mu$ ($\delta_\mu$ small) converge to their true Q-values;
error grows gracefully as actions move OOD.

\section{Experimental Findings}

\begin{itemize}[itemsep=2pt,parsep=0pt]
  \item \textbf{D4RL locomotion} (HalfCheetah, Hopper, Walker2D, Ant):
    BAC-PE improves over Diffusion-QL by 3--7 normalized score points on
    medium-expert datasets, the most challenging offline distribution.
  \item \textbf{Antmaze} (medium-play, large-diverse): 8--12 success rate
    improvement over CQL, which is excessively conservative on long-horizon tasks.
  \item \textbf{Robosuite manipulation} (Lift, PickPlace): 15\% improvement
    in task completion over TD3+BC, where multimodal action distributions are critical.
\end{itemize}

\section{Ablations and Interpretation}

\textbf{Advantage coefficient $\alpha$.} Setting $\alpha = 0$ (no correction) gives
standard diffusion Q-learning; increasing $\alpha$ improves performance up to $\alpha
\approx 1$ and then degrades, consistent with the theoretical bound: too large
$\alpha$ over-weights behavior-policy information at the expense of the true Q-values.

\textbf{Diffusion vs.\ Gaussian policy.} Replacing the diffusion policy with a
Gaussian policy degrades performance by 5--8 points on manipulation tasks where
action distributions are strongly multimodal, confirming that representation quality
matters alongside the BAC-PE correction.

\textbf{$Q^\mu$ estimation quality.} Using a poorly estimated $Q^\mu$ (trained on
a small data subset) degrades performance, confirming that BAC-PE's advantage
is contingent on reliable behavior-policy value estimation.

\section*{Reference}

Botao Dong, Longyang Huang, Ning Pang, Hongtian Chen.
\textbf{Diffusion Policy with Behavioral Advantage Correction for Offline
Reinforcement Learning}. arXiv:2608.02332, August 2026.
\url{https://arxiv.org/abs/2608.02332}

\end{document}
